\documentclass[11pt,a4paper]{article}
\usepackage[utf8]{inputenc}
\usepackage[T1]{fontenc}
\usepackage[french]{babel}
\usepackage{geometry}
\usepackage{booktabs}
\usepackage{hyperref}

\geometry{margin=2.2cm}

\title{LABeCO2 -- Documentation des donnees}
\author{LABeCO2}
\date{Mise a jour : 2026-05-18}

\begin{document}
\maketitle

\section{Source de verite}

LABeCO2 utilise SQLite comme source de verite pour les facteurs, les
consommables, les codes NACRES, les sources et les contributions.

\begin{itemize}
  \item \texttt{data/labeco2\_reference.sqlite} : base de reference livree avec l'application.
  \item \texttt{private/labeco2.sqlite} : base de travail locale en mode developpement.
  \item \texttt{scenarios/manips\_type.sqlite} : base des manipulations types.
\end{itemize}

Au lancement, l'application resout le chemin SQLite de travail. Si cette base
n'existe pas, elle est initialisee par copie de
\texttt{data/labeco2\_reference.sqlite}.

\section{Tables principales}

\begin{tabular}{lll}
\toprule
Table & Role & Remarques \\
\midrule
\texttt{purchase\_factors} & Facteurs achats/prix & Facteurs historiques GES/Labo1point5 consolides \\
\texttt{commercial\_products} & Produits commerciaux & Solides et liquides commerciaux \\
\texttt{product\_components} & Composants materiaux & Composants calculables avec masse renseignee \\
\texttt{materials} & Materiaux & Reliens aux facteurs d'emission materiaux \\
\texttt{emission\_factors} & Facteurs CO2e & Facteurs liquides et materiaux \\
\texttt{nacres\_codes} & Codes NACRES & Alimente les listes de selection de l'IHM \\
\texttt{transport\_factors} & Facteurs transport & Origine, distance et facteur par kg \\
\texttt{sources} & Sources documentaires & Traçabilite des ajouts et validations \\
\texttt{contributors} & Contributeurs & Auteurs des donnees \\
\bottomrule
\end{tabular}

\section{Validation et revisions}

Les tables metier portent un statut : \texttt{draft}, \texttt{validated} ou
\texttt{deprecated}. Lorsqu'une entree validee est modifiee, une revision est
creee et l'ancienne entree est conservee comme depreciee. L'outil
\texttt{tools/lab\_admin.py} permet de valider, fusionner et auditer les
contributions SQLite.

\section{Exports et contributions}

Les contributions peuvent etre exportees et importees via les outils :

\begin{itemize}
  \item \texttt{tools/export\_contribution.py}
  \item \texttt{tools/import\_contribution.py}
  \item \texttt{tools/export\_excel.py}
\end{itemize}

L'historique utilisateur de l'application est exporte en JSON, CSV ou Excel.
Les scenarios de calcul sont stockes hors des bases de reference afin de
separer les donnees metier validees des sessions utilisateur.

\section{Controles de coherence}

Les controles recommandes avant diffusion sont :

\begin{itemize}
  \item \texttt{pytest -q}
  \item \texttt{PRAGMA quick\_check}
  \item \texttt{PRAGMA foreign\_key\_check}
  \item audit via \texttt{tools/lab\_admin.py}
\end{itemize}

\end{document}
